\subsection{Требования по организации обмена данными с системами управления гофроагрегатом J.S.Machine}
\label{sec:BHS}

СИСТЕМА должна обмениваться данными с системами управления гофроагрегатом J.S.Machine. Программа управления заказами - ACS.

Обмен данными должен производиться из формы редактирования документа ''Выработка гофроагрегата'' (см. \ref{doc:Doc_ProdGA}).

\subsubsection{Настройка подключения для оборудования}

Для подключения к гофроагрегату J.S.Machine (Программа управления заказами - ACS)
в параметрах справочника ''Оборудование'' для гофроагрегата J.S.M. должны быть заполненные следующие данные.

\begin{enumerate}
\item Установлен флаг «Обмен с гофроагрегатом» – влияет на доступность команд обмена в документе выработки;
\item Интервал опроса БД, сек. – частота обновления данных из базы ACS;
\item Строка подключения с БД – параметры подключения к базе ACS, где:
\begin{itemize}
    \item DRIVER –- тип СУБД (MS SQL, MySQL и т.д.);
    \item SERVER –- адрес базы данных ACS в сети предприятия;
    \item ID –- имя пользователя (По умолчанию  для ACS ''sa'');
    \item PWD –- пароль (По умолчанию  для ACS ''sa'');
    \item DATABASE –- имя базы данных (По умолчанию для ACS ''[asc\_dev]'');
     
\end{itemize}
\end{enumerate}


\subsubsection{Структура файлов на JSACS}

% \setCJKmainfont{SimSun}
% \begin{CJK*}{UTF8}{zhsong}
Основные таблицы для работы

\point{Таблица Acs\_Order}

Таблица, содержащая список раскроев и их параметров.

% \begin{itemize}
%     \item fID – идентификатор раскроя (ID Задания);
%  \item 	Web – формат;
%  \item 	Remark – комментарий;
%  \item 	Trim – кромка;
%  \item 	Linedepth – глубина рилёвки;
%  \item 	SN – порядковый номер раскроя;
%  \item 	user\_fID – идентификатор пользователя;
%  \item 	AuditFields\_InsertDate – дата создания;
%  \item 	AuditFields\_UpdateDate – дата изменения;
%  \item 	LineType\_fID – идентификатор типа рилёвки (всегда равен 1);
%  \item 	Finished – признак, что раскрой был выполнен;
%  \item 	ChangSpd, BestSpd, maxSpd – параметры скорости работы;
%  \item 	Cust\_SN – порядковый номер (задаётся аналогично полю SN);
%  \item 	Cust\_SchNo – номер плана (не указывается);
%  \item 	Cust\_ProdNo – артикул (не указывается);
%  \item 	PaperMass\_fID – композиция;
%  \item 	PaperMassFlute\_fID – идентификатор профиля.
% \end{itemize}


% Table generated by Excel2LaTeX from sheet 'Раскрои'
\scriptsize
\begin{longtable}{|p{56mm}|p{56mm}|p{56mm}|}
\hline
{\bf \parbox[c][5mm]{56mm}{\centeringПоле}} & {\bf \parbox[c]{56mm}{\centeringТип}} & {\bf \parbox[c]{56mm}{\centeringПримечание}} \\
\hline
\parbox[c][5mm]{56mm}{Id} & \parbox{56mm}{Уникальный идентификатор} & \parbox{56mm}{Идентификатор Раскроя} \\
\hline
\parbox[c][5mm]{56mm}{Sn} & \parbox{56mm}{Число (4, 0)} & \parbox{56mm}{Порядковый номер Раскроя (10 - 9999)} \\
\hline
\parbox[c][5mm]{56mm}{OrderId} & \parbox{56mm}{Число (10, 0)} & \parbox{56mm}{Номер Раскроя} \\
\hline
\parbox[c][5mm]{56mm}{Web} & \parbox{56mm}{Число (4, 0)} & \parbox{56mm}{Формат} \\
\hline
\parbox[c][5mm]{56mm}{Trim} & \parbox{56mm}{Число (10, 4)} & \parbox{56mm}{Обрезь, мм} \\
\hline
\parbox[c][5mm]{56mm}{Trimmed} & \parbox{56mm}{Число (1, 0)} & \parbox{56mm}{Признак ''Без обрези''} \\
\hline
\parbox[c][5mm]{56mm}{Linedepth} & \parbox{56mm}{Число (10, 4)} & \parbox{56mm}{Глубина рилёвки, мм * 10} \\
\hline
\parbox[c][5mm]{56mm}{PaperMassGuid} & \parbox{56mm}{Уникальный идентификатор} & \parbox{56mm}{Идентификатор Композиции} \\
\hline
\parbox[c][5mm]{56mm}{PaperMassCode} & \parbox{56mm}{Строка (50)} & \parbox{56mm}{Наименование Композиции} \\
\hline
\parbox[c][5mm]{56mm}{PaperMassFluteGuid} & \parbox{56mm}{Уникальный идентификатор} & \parbox{56mm}{Идентификатор Профиля} \\
\hline
\parbox[c][5mm]{56mm}{FluteCode} & \parbox{56mm}{Строка (50)} & \parbox{56mm}{Наименование Профиля} \\
\hline
\parbox[c][5mm]{56mm}{LineTypeGuid} & \parbox{56mm}{Уникальный идентификатор} & \parbox{56mm}{Идентификатор Типа рилёвок} \\
\hline
\parbox[c][21mm]{56mm}{LineTypeValue} & \parbox{56mm}{Число (1, 0)} & \parbox{56mm}{Значение Типа рилёвок: 1 = ПапаМама; 2 = ПапаПлоский; 3 = ПапаПапа} \\
\hline
\parbox[c][5mm]{56mm}{LineTypeCode} & \parbox{56mm}{Строка (50)} & \parbox{56mm}{Наименование Типа рилёвок} \\
\hline
\parbox[c][5mm]{56mm}{Finished} & \parbox{56mm}{Число (1, 0)} & \parbox{56mm}{Признак завершённого Раскроя} \\
\hline
\parbox[c][5mm]{56mm}{ChangSpd} & \parbox{56mm}{Число (4, 0)} & \parbox{56mm}{Значение = 200} \\
\hline
\parbox[c][5mm]{56mm}{BestSpd} & \parbox{56mm}{Число (4, 0)} & \parbox{56mm}{Значение = 240} \\
\hline
\parbox[c][5mm]{56mm}{MaxSpd} & \parbox{56mm}{Число (4, 0)} & \parbox{56mm}{Значение = 200} \\
\hline
\parbox[c][5mm]{56mm}{PrePress} & \parbox{56mm}{Число (1, 0)} & \parbox{56mm}{Значение = 0} \\
\hline
\caption{Таблица Acs\_Order}\label{tab:Acs_Order}
\end{longtable}  
\normalsize


\point{Таблица 	Acs\_OrderItem}

Таблица, содержащая список заказов в раскроях.

% Table generated by Excel2LaTeX from sheet 'Заказы'
\scriptsize
\begin{longtable}{|p{56mm}|p{56mm}|p{56mm}|}
\hline
{\bf \parbox[c][5mm]{56mm}{\centeringПоле}} & {\bf \parbox[c]{56mm}{\centeringТип}} & {\bf \parbox[c]{56mm}{\centeringПримечание}} \\
\hline
\parbox[c][5mm]{56mm}{Id} & \parbox{56mm}{Уникальный идентификатор} & \parbox{56mm}{Идентификатор Заказа} \\
\hline
\parbox[c][5mm]{56mm}{OrderId} & \parbox{56mm}{Число (10, 0)} & \parbox{56mm}{Номер Раскроя} \\
\hline
\parbox[c][5mm]{56mm}{OrderItemId} & \parbox{56mm}{Число (10, 0)} & \parbox{56mm}{Номер Заказа} \\
\hline
\parbox[c][5mm]{56mm}{OrderGuid} & \parbox{56mm}{Уникальный идентификатор} & \parbox{56mm}{Идентификатор Раскроя} \\
\hline
\parbox[c][5mm]{56mm}{CustomerGuid} & \parbox{56mm}{Уникальный идентификатор} & \parbox{56mm}{Идентификатор Контрагента} \\
\hline
\parbox[c][5mm]{56mm}{CustomerID} & \parbox{56mm}{Строка (100)} & \parbox{56mm}{Код Контрагента} \\
\hline
\parbox[c][5mm]{56mm}{CustomerCode} & \parbox{56mm}{Строка (100)} & \parbox{56mm}{Наименование Контрагента} \\
\hline
\parbox[c][5mm]{56mm}{Sec} & \parbox{56mm}{Число (1, 0)} & \parbox{56mm}{Количество полос} \\
\hline
\parbox[c][5mm]{56mm}{StkModelGuid} & \parbox{56mm}{Уникальный идентификатор} & \parbox{56mm}{Идентификатор упаковки} \\
\hline
\parbox[c][5mm]{56mm}{StkModelCode} & \parbox{56mm}{Строка (100)} & \parbox{56mm}{Наименование упаковки} \\
\hline
\parbox[c][5mm]{56mm}{Length} & \parbox{56mm}{Число (10, 0)} & \parbox{56mm}{Длина заготовки, мм} \\
\hline
\parbox[c][5mm]{56mm}{Width} & \parbox{56mm}{Число (10, 4)} & \parbox{56mm}{Ширина заготовки, мм} \\
\hline
\parbox[c][5mm]{56mm}{Stk} & \parbox{56mm}{Число (10, 0)} & \parbox{56mm}{Количество в пачке} \\
\hline
\parbox[c][5mm]{56mm}{Qty} & \parbox{56mm}{Число (10, 0)} & \parbox{56mm}{Количество резов} \\
\hline
\parbox[c][10mm]{56mm}{StkNumber} & \parbox{56mm}{Число (4, 0)} & \parbox{56mm}{Количество пачек (по умолчанию = 20)} \\
\hline
\parbox[c][15mm]{56mm}{Level} & \parbox{56mm}{Число (1, 0)} & \parbox{56mm}{Стол: 0 = Верхний 1 = Нижний} \\
\hline
\parbox[c][5mm]{56mm}{IsNorMal} & \parbox{56mm}{Число (1, 0)} & \parbox{56mm}{Значение = 1} \\
\hline
\parbox[c][5mm]{56mm}{Lines} & \parbox{56mm}{Строка (100)} & \parbox{56mm}{Строка рилёвок} \\
\hline
\parbox[c][5mm]{56mm}{CustOrderId} & \parbox{56mm}{Число (1, 0)} & \parbox{56mm}{Значение = 1} \\
\hline
\parbox[c][5mm]{56mm}{ProductName} & \parbox{56mm}{Строка (100)} & \parbox{56mm}{Наименование изделия} \\
\hline
\parbox[c][5mm]{56mm}{PlanSheet} & \parbox{56mm}{Строка (100)} & \parbox{56mm}{Номер Плана} \\
\hline
\parbox[c][5mm]{56mm}{PlanFinishedDate} & \parbox{56mm}{Дата} & \parbox{56mm}{Плановая дата завершения} \\
\hline
\parbox[c][5mm]{56mm}{IsPaperBox} & \parbox{56mm}{Число (1, 0)} & \parbox{56mm}{Значение = 0} \\
\hline
\parbox[c][5mm]{56mm}{IsQuickOut} & \parbox{56mm}{Число (1, 0)} & \parbox{56mm}{Значение = 0} \\
\hline
\parbox[c][5mm]{56mm}{Sn} & \parbox{56mm}{Число (4, 0)} & \parbox{56mm}{Порядковый номер Раскроя} \\
\hline
\parbox[c][5mm]{56mm}{OutPaperPath} & \parbox{56mm}{Число (1, 0)} & \parbox{56mm}{Значение = 0} \\
\hline
\parbox[c][5mm]{56mm}{Trim} & \parbox{56mm}{Число (10, 4)} & \parbox{56mm}{Обрезь, мм} \\
\hline
\parbox[c][5mm]{56mm}{Web} & \parbox{56mm}{Число (4, 0)} & \parbox{56mm}{Формат} \\
\hline
\parbox[c][5mm]{56mm}{StkModelValue} & \parbox{56mm}{Число (1, 0)} & \parbox{56mm}{Значение = 0} \\
\hline
\parbox[c][5mm]{56mm}{OrderStaus} & \parbox{56mm}{Число (1, 0)} & \parbox{56mm}{Значение = 0?} \\
\hline
\caption{Таблица Acs\_OrderItem}\label{tab:Acs_OrderItem}
\end{longtable}  
\normalsize




\point{Таблица 	Acs\_CompleteOrder	}

Таблица, содержащая информацию о выполненных раскроях.

% Table generated by Excel2LaTeX from sheet 'Данные выработки'
\scriptsize
\begin{longtable}{|p{56mm}|p{56mm}|p{56mm}|}
\hline
{\bf \parbox[c][5mm]{56mm}{\centeringПоле}} & {\bf \parbox[c]{56mm}{\centeringТип}} & {\bf \parbox[c]{56mm}{\centeringПримечание}} \\
\hline
\parbox[c][5mm]{56mm}{Sn} & \parbox{56mm}{Число (4, 0)} & \parbox{56mm}{Номер раскроя на ГА} \\
\hline
\parbox[c][5mm]{56mm}{OrderGuid} & \parbox{56mm}{Уникальный идентификатор} & \parbox{56mm}{УИД строки Раскроя} \\
\hline
\parbox[c][5mm]{56mm}{OrderItemGuid} & \parbox{56mm}{Уникальный идентификатор} & \parbox{56mm}{УИД строки Заказа} \\
\hline
\parbox[c][5mm]{56mm}{StartTime} & \parbox{56mm}{Дата} & \parbox{56mm}{Дата начала выработки} \\
\hline
\parbox[c][5mm]{56mm}{EndTime} & \parbox{56mm}{Дата} & \parbox{56mm}{Дата окончания выработки} \\
\hline
\parbox[c][10mm]{56mm}{ProdQty} & \parbox{56mm}{Число (10, 0)} & \parbox{56mm}{Количество выпущенных заготовок, шт.} \\
\hline
\parbox[c][10mm]{56mm}{WastQty} & \parbox{56mm}{Число (10, 0)} & \parbox{56mm}{Количество бракованных заготовок, шт.} \\
\hline
\parbox[c][26mm]{56mm}{StopTime} & \parbox{56mm}{Число (10, 0)} & \parbox{56mm}{Общее время простоя, сек.
(Если значение <> 0, то в Журнал работы оборудования будет добавлена запись о простое)} \\
\hline
\caption{Таблица Acs\_CompleteOrder} \label{tab:Acs_CompleteOrder}
\end{longtable}  
\normalsize

 
 
%В данной таблице всегда отображается половина выполняемого раскроя в текущий момент. Например, если раскрой состоит из двух заказов, здесь будет строчка второго заказа с оперативно обновляемым фактом выработки, в то время как строки первого заказа нет. Однако, как только раскрой завершится, факт выработки по обоим заказам отразится корректно. Вероятно, эта особенность является недоработкой со стороны ACS.

\point{Таблица 	Acs\_ProdPresses	} 

Таблица, содержащая информацию о текущих данных выработки.

% ВНИМАНИЕ. Гофроагрегат J.S.Machine и система DCS записывает данные в указанную таблицу по раскроям ТОЛЬКО ПО ФАКТУ ОКОНЧАНИЯ ВЫПОЛНЕНИЯ РАСКРОЯ. 
% Промежуточная информация по раскрою НЕДОСТУПНА.

% Table generated by Excel2LaTeX from sheet 'Текущие данные'
\scriptsize
\begin{longtable}{|p{56mm}|p{56mm}|p{56mm}|}
\hline
{\bf \parbox[c][5mm]{56mm}{\centeringПоле}} & {\bf \parbox[c]{56mm}{\centeringТип}} & {\bf \parbox[c]{56mm}{\centeringПримечание}} \\
\hline
\parbox[c][5mm]{56mm}{OrderItemGuid} & \parbox{56mm}{Уникальный идентификатор} & \parbox{56mm}{УИД строки Заказа} \\
\hline
\parbox[c][5mm]{56mm}{StartTime} & \parbox{56mm}{Дата} & \parbox{56mm}{Дата начала выработки} \\
\hline
\parbox[c][5mm]{56mm}{EndTime} & \parbox{56mm}{Дата} & \parbox{56mm}{Дата окончания выработки} \\
\hline
\parbox[c][10mm]{56mm}{ProdArea} & \parbox{56mm}{Число (18, 4)} & \parbox{56mm}{Площадь выпущенных заготовок, м2} \\
\hline
\parbox[c][10mm]{56mm}{BadArea} & \parbox{56mm}{Число (18, 4)} & \parbox{56mm}{Площадь бракованных заготовок, м2} \\
\hline
\caption{Таблица Acs\_ProdPresses}\label{tab:Acs_ProdPresses}
\end{longtable}  
\normalsize




Данная таблица всегда содержит N+1 строк, где N – количество заказов в раскрое. Последняя строка отражает суммарные данные по смене.


% \end{CJK*}

% Для выгрузки в память гофроагрегата СИСТЕМА выгружает файл с производственными заказами и раскроями в общую папку компьютера гофроагрегата. 
% Имя общей папки --- ''IPAddress/RcvOrder'' или ''ComputerName/RcvOrder'',
% где IPAddress  -- IP-адрес компьютера с утсновленонй системой DCS, ComputerName -- сетевое имя компьютера DCS.
% Папка RcvOrder должна находится в корневом каталоге диска D компьютера DCS.

%  СИСТЕМА должна иметь доступ к общей папке. Кодировка файлов --- ASCII. Каждая строка в файлах должна заканчиваться стандартной комбинацией символов конца строки (0x0D, 0x0A). Строка состоит из последовательности полей. Отдельное поле представляет собой комбинацию символов. Используемые форматы полей представлены в таблице \ref{tab:jlm_structure_of_output}.

% DCS считывает и удаляет текст OfficeH.txt

%\scriptsize
%\pc
%\begin{longtable}{|p{20mm}|p{40mm}|p{60mm}|}
%\hline
%{\bf Обозначение} & {\bf Тип значения} & {\bf } \endhead
%\hline AN & буквенно-цифровой & значение выравнено по левому краю\\
%\hline N & числовой & значение выравнено по правому краю\\
%\hline N/B & числовой или пустой & значение выравнено по правому краю\\
%\hline D & дата (ГГГГММДД) & значение состоит из 8 цифр\\
%\hline T & время (ЧЧММСС) & значение состоит из 6 цифр\\
%\hline ID & идентификатор установки & значение состоит из 20 знаков\\
%\hline PID & идентификатор бумаги & значение состоит из 44 знаков\\
%\hline
%\caption{Форматы полей, используемые в файлах обмена с гофроагрегатами BHS}
%\label{tab:bhs_formats_of_fields}
%\end{longtable} 
%
%\normalsize
%
%Идентификатор установки всегда состоит из 20 знаков. Структура идентификатора установки представлена в таблице \ref{tab:bhs_structure_of_id}.
%
%\scriptsize
%\pc
%\begin{longtable}{|p{20mm}|p{80mm}|p{20mm}|}
%\hline
%{\bf Длина} & {\bf Значение} & {\bf Формат} \endhead
%\hline 16 & идентификатор установки & AN\\
%\hline 3 & серийный номер & N/B\\
%\hline 1 & идентификатор этапа работы гофроагрегата по заданной установке & N/B\\
%\hline
%\caption{Структура идентификатора установки}
%\label{tab:bhs_structure_of_id}
%\end{longtable}
%
%\normalsize
%
%Идентификатор этапа работы гофроагрегата по заданной установке используется в ситуациях, когда работа гофроагрегата по заданной установке разделяется на этапы. Например, работа гофроагрегата по заданной установке останавливается вечером и возобновляется утром. В такой ситуации идентификатор уже выполненного этапа работы гофроагрегата может быть равен 1, а идентификатор еще невыполненного этапа –-- 2.
%
%Идентификатор бумаги всегда состоит из 44 знаков. Структура идентификатора бумаги представлена в таблице \ref{tab:bhs_structure_of_pid}.
%
%\scriptsize
%\pc
%\begin{longtable}{|p{20mm}|p{80mm}|p{20mm}|}
%\hline
%{\bf Длина} & {\bf Значение} & {\bf Формат} \endhead
%\hline 20 & код бумаги & AN\\
%\hline 4 & плотность бумаги (г/м$^2$) & N\\
%\hline 20 & наименование бумаги (дополнительные коды) & AN/B\\
%\hline
%\caption{Структура идентификатора бумаги}
%\label{tab:bhs_structure_of_pid}
%\end{longtable}
%
%\normalsize



\subsubsection{Функциональные требования}

\point{Выгрузить раскрои в J.S.Machine}

Выгружать СИСТЕМА должна раскрои без признака ''Выгружен''. В форме редактирования документа ''Выработка гофроагрегата'' в табличной части ''Раскрои'' СИСТЕМА должна позволять выделять строки (несколько строк) раскроев. 
По нажатию на кнопку ''Выгрузить раскрои'' СИСТЕМА должна 
автоматически подключаться к базе данных гофроагрега и выгрузить отмеченные раскрои. У выгруженных раскроев заполнится поле «Номер связки» – это порядковый номер раскроя в очереди гофроагрегата. 


Выгруженные раскрои должны быть помечены в таблице  признаком ''Выгружен''. 

В течение смены новые раскрои могут также добавляться и выгружаться без необходимости остановки чтения данных гофроагрегата.


Для формирования запроса требуется следующая структура данных:
\pc
\begin{longtable}{|p{4cm}|p{4cm}|p{8cm}|}
\hline
{\bf Наименование} & {\bf Тип данных} &  {\bf Комментарий} \endhead
    \hline
    УИД раскроя &	Уникальный идентификатор	& Уникальный идентификатор Раскроя \\
    \hline
    Номер раскроя &	Строка (50)	& ID Задания \\
    \hline
    Композиция &	Строка (50)&	Композиция раскроя, состоящая из кодов слоёв
(см.  "Общие правила формирования заданий", п.4)\\
    \hline
    Профиль &	Строка (50)&	 \\
    \hline
    Формат &	Число (4, 0)	 \\
        \hline
     Обрезь&	Число (4, 0)&	Если равна 0, то к формату прибавляется 1, 
в поле Trim ставится 1, в поле Trimmed ставится 0\\
        \hline
    Тип рилёвок&	Число (4, 0)	& Возможные значения: 
1 – Папа-Мама, 2 – Папа-Папа, 3 – Прочее\\
        \hline
    Комментарий&	Строка (255)&	Дублируется ID Задания для оператора\\
        \hline
    УИД заказа&	Строка (50)& Уникальный идентификатор Заказа \\
        \hline
    Номер заказа&	Строка (50)&	 \\
        \hline
    Стол&	Строка (10)& 0 - Верхний, 1 - Нижний \\
        \hline
     Количество полос&	Число (4, 0)&	 \\
        \hline
     Рилёвки с 1 по 5&	Строка (100)&	Если нет рилёвок, ставится ширина заготовки\\
        \hline
     Длина заготовки&	Число (4, 0)	 &\\
     \hline
     Ширина заготовки&	Число (4, 0)	 &\\
    \hline
     Кол-во резов	&Число (4, 0)&	 \\
    \hline
     Кол-во пачек	&Число (4, 0)&	По умолчанию = 20\\
    \hline
     Кол-во в пачке&	Число (4, 0)&\\	%(см. "Общие правила формирования заданий", п. 5)
    \hline
     УИД контрагента & Уникальный идентификатор &  Уникальный идентификатор контрагента\\
    \hline
     Наименование контрагента&	Строка (50)&	Для создания нового, если по идентификатору не нашли\\
    \hline
     Комментарий &	Строка (100)&	ID Задания\\
    \hline
\end{longtable}

Запрос должен состоять из следующих подзапросов:
\begin{enumerate}
    \item 	Объявление начала транзакции;
    \item 	Создание и заполнение параметров работы оборудования:
    \begin{enumerate}
    \item 	Создание новых слоёв (Acs\_Paper);
    \item 	Создание новых профилей (Acs\_Flute);
    \item 	Создание новых связок ''Композиция + Профиль'' (Acs\_PaperMass);
    \item 	Создание параметров работы для связок ''Композиция + Профиль'' (Acs\_PaperReserParameters);
    \item 	Создание параметров глубины рилёвок для связок ''Профиль + Тип рилевки'' (Acs\_FlutesLineDepth);
    \end{enumerate}
    \item 	Создание/обновление контрагентов (Acs\_Customer);
    \item 	Создание нового раскроя (Acs\_Order);
    \item 	Создание состава раскроя (Acs\_OrderItem);
    \item 	Уничтожение временных таблиц, объявление завершения транзакции, откат изменений при наличии ошибок.
\end{enumerate}



%Для проверки того, что раскрои встали в очередь компьютера гофроагрегата, необходимо считать файл ''SEQUENCE.DAT'', структура строки которого представлена в таблице \ref{tab:bhs_structure_of_row_of_sequencefile}, и проверить, содержит ли он строки, соответствующие выгруженным раскроям. Если строка с идентификатор установки, соответствующей выгруженному раскрою, присутствует в файле и не содержит ошибки, тогда в табличной части для него необходимо установить признак ''Выгружен''. 
%Если же такая строка отсутствует в файле, либо содержит ошибку, значит раскрой не был корректно выгружен и необходимо оповестить пользователя об этом.

\bigskip

\point{Начать чтение данных с ГА}
\label{sec:startreading}

Новая кнопка и команда ''Начать чтение данных с ГА''  должна быть добавлена в форму редактирования документа ''Выработка гофроагрегата'' (см. \ref{doc:Doc_ProdGA}).

При считывании данных СИСТЕМА должна выполнять считывание данных с гофроагрегата с интервалом ''Интервал опроса БД'', значение которого установлено в свойствах Оборудования (Справочник ''Оборудования'').
При считывании должен вызываться метод ''Считать данные с ГА'' (см. \ref{poЧисло (4, 0):readdatafromcorrugator}).

\bigskip

\point{Закончить чтение данных с ГА}
\label{sec:stopreading}

Новая кнопка и команда ''Закончить чтение данных с ГА'' должна быть добавлена в форму редактирования документа ''Выработка гофроагрегата'' (см. \ref{doc:Doc_ProdGA}).

По команде оператора необходимо прекратить опрос базы данных гофроагрегата.

Документ сохранить автоматически.


\bigskip

\point{Считать данные с ГА}
\label{poЧисло (4, 0):readdatafromcorrugator}


Общие правила для формирования заданий:

\begin{enumerate}

    \item Уникальные идентификаторы Раскроя и Заказов формируются автоматически при добавлении Раскроя в документ выработки;
    \item Номер раскроя формируется из Номера документа и порядкового Номера раскроя в таблице, который затем будет преобразован в число (например, номер документа - ''000007110'', номер раскроя - ''2'' --> ''71102'');
    \item Номер заказа будет преобразован в число (например, ''4378.3'' --> ''43783'');

    \item Композиция раскроя формируется на основе кодов ACS из слоёв картона, от 1 до 7 слоя. При этом, если слоя нет, то вместо него ставятся символы ''-{}-''. 
    \begin{itemize}
      \item 	Например, раскрой состоит из трёх слоёв: Братск-115, Омск-125, Алтай-125;
       \item 	У каждого слоя определён код: Братск-115 = B3, Омск-125 = O4, Алтай-125 = A3;
       \item 	Следовательно, композиция для задания примет вид: ''B3O4A3-{}-{}-{}-{}-{}-{}-{}-''.
    \end{itemize}
    
    \item Если в заказе всего 1 заказ, и для раскроя установлен признак ''Распределять по столам'':
    \begin{itemize}
	\itemЕсли количество полос больше 2, тогда заказ дробится пополам по количеству полос, причём на нижний стол должна пойти часть с большим количеством полос (например, если изначально в задании 5 полос, то 2 полосы пойдут на верхний стол, а 3 - на нижний).

	\itemЕсли количество полос меньше или равно 2, то такой заказ встанет на нижний стол. Но к нему на верхний стол будет добавлен фиктивный заказ – копия первого с нулевым количеством, длиной/шириной заготовок, полос и т.д.
    \end{itemize}
    \item 	Если обрезь получается нулевой, то к формату добавляется единица, в поле Trim устанавливается 1, а в поле Trimmed устанавливается 0.
    \item 	Количество в пачке (поле Stk) определяется по правилу:
        \begin{itemize}
       \item 	Для гофролистов количество в пачке равно соответствующему полю технологической карты, если оно заполнено;
       \item 	Для профиля BC: Ширина заготовки <= 300: 150; >300: 250;
       \item 	Для профиля BE: Ширина заготовки <= 300: 200; >300: 300;
       \item 	Для профиля C: Ширина заготовки <= 300: 200; >400: 350; >300: 400;
       \item 	Для профиля B: Ширина заготовки <= 300: 300; >300: 550;
       \item 	Для профиля E: Ширина заготовки <= 300: 400; >300: 750;
       \item 	Значение по умолчанию – 300 шт./пач.
    \end{itemize}
    
\end{enumerate}


% Получение текущих данных выработки (Acs\_ProdPresses):

% \pc
% \begin{longtable}{|p{4cm}|p{4cm}|p{8cm}|}
% \hline
% {\bf Наименование} & {\bf Тип данных} & {\bf Комментарий} \endhead
% \hline
% OrderItemGuid & Уникальный идентификатор & Уникальный идентификатор Заказа \\
% \hline
% StartTime & Дата & Дата начала выработки \\
% \hline
% EndTime & Дата & Дата окончания выработки \\
% \hline
% ProdArea & Число (18, 4) & Площадь выпущенных заготовок, м2 \\
% \hline
% BadArea & Число (18, 4) & Площадь бракованных заготовок, м2 \\
% \hline
% \end{longtable}

% Получение данных завершённой выработки (Acs\_CompleteOrder):

% \pc
% \begin{longtable}{|p{4cm}|p{4cm}|p{8cm}|}
% \hline
% {\bf Наименование} & {\bf Тип данных} & {\bf Комментарий} \endhead
% \hline
% Sn & Число (4, 0) & Порядковый номер раскроя \\
% \hline
% OrderGuid & Уникальный идентификатор & Уникальный идентификатор Раскроя \\
% \hline
% OrderItemGuid & Уникальный идентификатор & Уникальный идентификатор Заказа \\
% \hline
% StartTime & Дата & Дата начала выработки \\
% \hline
% EndTime & Дата & Дата окончания выработки \\
% \hline
% ProdQty & Число (4, 0) & Количество выпущенных заготовок, шт. \\
% \hline
% WastQty & Число (4, 0) & Количество бракованных заготовок, шт. \\
% \hline
% StopTime & Число (4, 0) & Общее время простоя, сек. Если значение <> 0, то в Журнал работы оборудования будет добавлена запись о простое \\
% \hline
% \end{longtable}